% Source: https://artofproblemsolving.com/community/c3h1848454_some_problems
% Stable identity: aops-topic-1848454-q02
\textbf{(a)} The part under the first ``$2+$'' is $\sqrt{3+x}$, so $x=\sqrt{2+\sqrt{3+x}}$. Squaring, $x^2-2=\sqrt{3+x}$, and squaring again,
\[
\left(x^2-2\right)^2=3+x\iff\boxed{x^4-4x^2-x+1=0}.
\]

\textbf{(b)} Let $g(t)=t^4-4t^2-t+1$. Since $x^2-2=\sqrt{3+x}>0$, we have $x>\sqrt2$. On $[\sqrt2,\infty)$ the function $g$ is convex, because $g''(t)=12t^2-8>0$ there, and $g(\sqrt2)=4-8-\sqrt2+1<0$. A convex function that is negative at the left end of an interval crosses zero at most once to the right of it, and $g(t)\to\infty$, so $g$ has exactly one root greater than $\sqrt2$, and it must be $x$. Finally
\[
g(2.06)\approx-0.026<0,\qquad g(2.07)\approx0.151>0,
\]
so $\boxed{2.06<x<2.07}$ (in fact $x\approx2.0615$).

(The quartic has no rational roots and does not factor over the rationals, so there is no simpler closed form to ``simplify'' to; the value is best described as this root.)
